% ===== 02-background.tex =====
\section{Background}
\label{sec:background}
\textbf{Agent executions.} An agent runtime (harness) drives a loop in which a language model selects tools; each invocation is a \emph{tool call}, and the steps devoted to one goal form a \emph{task}~\cite{yao2023react}. In this paper, an \emph{execution} is the semantic unit to which network activity should be attributed. Its correspondence to operating-system entities depends on how the runtime organizes work: one execution may spawn many processes, and one worker process may serve many executions.

\textbf{Kernel identity signals.} A PID identifies the task that issues a system call. \emph{Process lineage} associates descendants with the root of an execution tree; it discriminates executions only while they live in distinct trees. \emph{Cgroups} group processes for accounting and resource control; their membership discriminates executions only when the group layout separates them~\cite{zheng2026agentcgroup}. \emph{Application context}, such as session, task, or tool-call identifiers, can be bound to runtime facts~\cite{byteyellow2026agentprovenance,otel2026genai}. It can represent distinctions absent from the process structure, but its correctness depends on the binding.

\textbf{Attribution outcomes.} We classify the resolver output for each observed event against a reference origin: \emph{correct} (C) when the single selected origin matches the reference; \emph{ambiguous} (A) when the output is a set of candidates; \emph{missing} (M) when no origin is assigned; and \emph{incorrect} (I) when the single selected origin differs from the reference. For management, I is the costly outcome: it credits traffic to the wrong execution without warning. A, by contrast, is visible and can be handled by policy.
